% !TEX root = ../../tree_notes.tex
% =====================================================================
%  6.4  Recover the Structure by Queries / Interactive   sec:c25-query
%  Exports: def:c25-query  prop:c25-halve  thm:c25-logn  lst:c25-guess
%           ex:c25-query-run  prob:c25-guess  lst:c25-guess  (practice: LC 375)
% =====================================================================
\section{Recover the Structure by Queries (Interactive)}\label{sec:c25-query}

\begin{previouslybox}
\begin{enumerate}[nosep, leftmargin=1.6em]
  \item What does a single comparison in binary search do to the size of the candidate set? Why
        $\log_2 n$ comparisons in total? (\chref{c01}, \chref{c24} preview)
  \item In \secref{c25-parent}, what $\bigO{1}$ fact about the root made it readable from an array?
        (\propref{c25-root})
  \item How could a \emph{distance} or \emph{ancestor} answer about a tree be produced by the
        timestamp machinery? (\chref{c24}, \thmref{c24-anc})
  \item If a yes/no answer is truthful and splits the world in half, how many do you need to pin
        one of $N$ possibilities?
\end{enumerate}
\end{previouslybox}

\begin{topicintro}
Every topic so far \emph{handed} you the data. Here the data is hidden behind an \term{oracle}:
you may only \emph{ask questions}, and each answer is a small, costly bit of information you must
\emph{use adaptively} -- your next query depends on the previous answers. Reconstructing a hidden
object by queries is the same skill as the parent/degree/identification problems of
\secref{c25-parent}--\secref{c25-identify}, turned inside out: instead of reading a full array you
\emph{elicit} it, one strategic question at a time. The purest instance -- and our worked anchor
LC 374 -- is \term{guess the number}: pick from $[1,n]$, an oracle replies ``higher / lower /
correct'', and the optimal strategy is textbook binary search, which provably uses the minimum
possible number of questions (\thmref{c25-logn}). Two ideas carry the whole section: a truthful
answer should \emph{halve} the candidate space, and the number of candidates you can distinguish is
$2^{\text{\#answers}}$, so you cannot beat $\lceil\log_2\rceil$. Where the \emph{cost} of a wrong
guess matters (LC 375) the halving point shifts from the middle to a min-max choice -- an interval
DP -- and where the hidden object is a \emph{tree} rather than a number, the same adaptive idea
becomes group-testing on \term{distance}/\term{ancestor} queries answered by the Chapter~5 oracle
\chref{c24}. This is construction from data where you get to \emph{choose} the data.
\textbf{Prerequisites.} Binary-search intuition; the Chapter~5 distance/ancestor oracle
(\chref{c24}) used as a black box. This section adds the \emph{query model, the
information-theoretic lower bound, and adaptive vs.\ cost-aware strategies}.

\textbf{After this section you can}
\begin{itemize}[nosep]
  \item model a problem as an adaptive query protocol and identify the candidate space;        \hfill\lv{2}
  \item apply the halving principle and prove $\lceil\log_2 N\rceil$ questions are necessary;   \hfill\lv{2}--\lv{3}
  \item implement optimal search against a three-way (higher/lower/equal) oracle (LC 374);      \hfill\lv{2}
  \item recognise when the optimal split is \emph{not} the middle (worst-case cost, LC 375).    \hfill\lv{3}--\lv{4}
\end{itemize}
\end{topicintro}

\begin{recognitionbox}[Pattern card -- reconstruct by asking the oracle]
\reckey{you do not see the hidden value/object; instead you are given a \emph{feedback function}
(a \code{guess()}, an API that answers ``higher/lower'', a distance/ancestor query, a yes/no test)
and must find the answer with as few \emph{queries} as possible, adapting each to the last.}
\recuse{
  \item ``guess the number / word / root'' with a comparison oracle -- binary search (\probref{c25-guess});
  \item ``find $x$ using $f(x)$ queries, minimise the \emph{number} of queries'' -- halve the space;
  \item ``guarantee a win, but a wrong guess \emph{costs} you'' -- min-max, not midpoint (LC 375);
  \item ``rebuild a hidden tree using distance / ancestor queries'' -- adaptive group testing on
        the \chref{c24} oracle.
}
\recnotuse{
  \item the full array/matrix is visible up front -- \secref{c25-parent}/\secref{c25-degree};
  \item you just need to search a \emph{sorted} array once, not drive an interactive protocol
        -- ordinary binary search, \chref{c01};
  \item answers may be \emph{lies} (Ulam's game) -- that needs error-correcting splits, out of scope.
}
One-line test: \emph{can you only learn the answer by \textbf{asking}, and does each reply shrink
the possibilities? Choose the question that halves -- this section.}
\end{recognitionbox}

\subsection{Intuition}

You are on a number line $[1,n]$; the secret $x$ is somewhere on it and a referee only says
``higher'', ``lower'' or ``got it'' after each of your guesses. The clever move is always to guess
the \emph{middle} of the still-possible interval: whatever the referee says, you discard half the
remaining candidates. The interval of possibilities shrinks $n\to n/2\to n/4\to\cdots$, so about
$\log_2 n$ guesses finish it -- and you cannot do better, because each three-way answer can at best
divide the world into three and the count of distinguishable secrets grows only exponentially in the
number of questions. This is the whole game: \emph{a good query is one whose answers leave you with
equally small sub-cases}.

\begin{figure}[H]
\centering
\begin{tikzpicture}[diagram, x=11mm, y=8mm]
  % candidate interval rows, shrinking
  \node[annot, anchor=east] at (0,0) {\small $[1,\,10]$};
  \foreach \k in {1,...,10} { \node[cellnode, fill=primaryblue!10] at (\k,0) {\k}; }
  \node[annot, anchor=east] at (0,-1.1) {\small ask $5$};
  \foreach \k in {1,...,10} { \node[cellnode, fill=primaryblue!6] at (\k,-1.1) {\k}; }
  \fill[difforange!25] (0.62,-1.48) rectangle (5.38,-0.72);   % discard 1..5 (guess=5 correct-branch shown)
  \node[annot, anchor=east] at (0,-2.2) {\small ans.\! high};
  \foreach \k in {1,...,10} { \node[cellnode, fill=primaryblue!6] at (\k,-2.2) {\k}; }
  \fill[successgreen!28] (5.62,-2.58) rectangle (10.38,-1.82); % keep 6..10
  \node[annot, align=left, font=\scriptsize] at (5.5,-3.5)
    {guess the midpoint $5$; if told ``secret is higher'' the survivors $[6,10]$ are half\\
     of what remains -- each truthful answer throws away $\approx$ half the candidates};
\end{tikzpicture}
\caption{Adaptive halving on a number line. One well-chosen question (the midpoint) turns a size-$10$
candidate set into a size-$\le5$ one no matter which branch the oracle confirms. Repeating gives a
decision \emph{tree} of depth $\lceil\log_2 n\rceil$ -- exactly the shape the tree chapters draw, but
here its nodes are \emph{questions}, not data.}\label{fig:c25-halve}
\end{figure}

\subsection{Definitions and Notation}

\begin{definition}{Query (oracle) model}{c25-query}
A \term{query model} is a tuple $(\mathcal{X},\mathcal{Q},\mathit{ans})$: a hidden target
$x\in\mathcal{X}$ (the candidate space), a set of allowed \term{queries} $q\in\mathcal{Q}$, and an
answer function $\mathit{ans}(q,x)$ returning a feedback symbol. An \term{adaptive strategy} chooses
each next query as a function of the answers so far, and must halt with a unique $x$. Its
\term{cost} is the number (or weighted total) of queries in the worst case over $x$.
\end{definition}

\noindent\textbf{In plain words.} You can't see $x$; you may only poke questions and get terse
truthful replies; a strategy is good if it always names $x$ using few questions, and \emph{adaptive}
means your next question depends on the last answer.

\begin{proposition}{Halving and the decision-tree bound}{c25-halve}
If every query has at most $b$ possible answers, then $k$ queries can distinguish at most $b^{k}$
targets, so any correct adaptive strategy needs
$\displaystyle k\ \ge\ \lceil \log_{b} |\mathcal{X}| \rceil$ queries in the worst case. For a
three-way comparison oracle ($b=3$: higher/lower/equal) on $N$ ordered secrets this is
$\lceil\log_3 N\rceil$ \emph{in principle}, and binary search attains the (slightly stronger,
because answers must respect order) bound of $\lceil\log_2(N{+}1)\rceil$ guesses.
\end{proposition}
\begin{proofbox}
Each answer can at best multiply the number of distinguishable outcomes by $b$: after $k$ queries the
transcript of answers has $\le b^{k}$ distinct values, and a correct strategy must map different
targets to different transcripts, so $|\mathcal{X}|\le b^{k}$, i.e.\ $k\ge\log_b|\mathcal{X}|$. The
decision process is literally a $b$-ary \emph{tree} (nodes $=$ queries, edges $=$ answers, leaves $=$
conclusions) of depth $k$ and hence $\le b^{k}$ leaves. For an ordered search each comparison splits
the remaining interval into a left block, a point, and a right block; always splitting so the larger
side is as small as possible yields the recurrence $T(N)=1+T(\lfloor N/2\rfloor)$ whose solution is
$\lceil\log_2(N{+}1)\rceil$, attained by guessing the midpoint. \end{proofbox}

\begin{theorem}{Binary search is query-optimal}{c25-logn}
Against a truthful higher/lower/equal oracle, guessing the midpoint of the current candidate interval
minimises the worst-case number of queries; it finds any secret in $[1,n]$ within
$\lceil\log_2(n{+}1)\rceil$ guesses, matching \propref{c25-halve}'s lower bound up to the base. No
strategy can guarantee fewer in the worst case.
\end{theorem}
\begin{insightbox}
Optimality is \emph{about the split}, not the midpoint per se: you want the answer that leaves the
\emph{largest} remaining branch as small as possible. For a pure \emph{count} of questions on a
uniformly ordered line, that is the midpoint -- binary search. \emph{When answers cost something}
(LC 375: you pay the number you guessed) the split is no longer the middle -- you minimise the worst
case \emph{sum}, giving $k+\max(\text{left},\text{right})$ rather than $\max$ over sizes. Same
decision tree, different objective.
\end{insightbox}

\subsection{The algorithm: optimal play against a three-way oracle}

\begin{lstlisting}[caption={Guess Number Higher or Lower -- overflow-safe binary search on a comparison oracle}, label={lst:c25-guess}]
#include <bits/stdc++.h>
using namespace std;

// Forward-declared oracle (provided by the judge):
//   returns -1 if num is HIGHER than the secret ("too high"),
//            1 if num is LOWER than the secret ("too low"),
//            0 if num equals the secret.
int guess(int num);

class Solution {
public:
    int guessNumber(int n) {
        int lo = 1, hi = n;
        while (lo <= hi) {
            int mid = lo + (hi - lo) / 2;      // NOT (lo+hi)/2 -- lo+hi overflows int (n up to 2^31-1)
            int g = guess(mid);
            if (g == 0)  return mid;           // hit
            if (g == -1) hi = mid - 1;         // mid too high -> secret is in [lo, mid-1]
            else         lo = mid + 1;         // mid too low  -> secret is in [mid+1, hi]
        }
        return -1;                             // unreachable for a valid secret in [1,n]
    }
};
\end{lstlisting}
\begin{complexitybox}
Time $\bigO{\log n}$ guesses -- the interval halves every iteration (\thmref{c25-logn}); at
$n=2^{31}{-}1$ that is $\le 31$ oracle calls. Memory $\bigO{1}$. The \emph{overflow guard}
\code{lo + (hi-lo)/2} is the real content of an ``easy'' problem: with $n$ up to $2^{31}{-}1$, the
naive \code{(lo+hi)/2} exceeds \code{INT\_MAX} and mis-splits.
\end{complexitybox}

\subsection{Worked Example}

\begin{example}{Finding $x=2$ in $[1,10]$ by adaptive halving}{c25-query-run}
Secret $x=2$. Start $[\mathit{lo},\mathit{hi}]=[1,10]$.
Guess $\mathit{mid}=1{+}(10{-}1)/2=5$; $\mathit{guess}(5)=-1$ (too high) $\Rightarrow$
$\mathit{hi}=4$, interval $[1,4]$. Guess $\mathit{mid}=2$; $\mathit{guess}(2)=-1$ (too high)
$\Rightarrow$ $\mathit{hi}=1$, interval $[1,1]$. Guess $\mathit{mid}=1$; $\mathit{guess}(1)=1$
(too low) $\Rightarrow$ $\mathit{lo}=2$, and the loop's $\mathit{mid}=2$ gives $0$ after at most
$\lceil\log_2 11\rceil=4$ guesses. Every step threw away about half the remaining candidates -- the
\figref{c25-halve} shrink in action, finishing within the \thmref{c25-logn} bound.
\end{example}
\begin{tracebox}[Why the midpoint is the right question]
From a size-$m$ interval a midpoint guess leaves $\lfloor m/2\rfloor$ candidates on the ``too high''
branch and $\lceil m/2\rceil-1$ on the ``too low'' branch, so the larger survivor is $\le\lceil m/2\rceil$.
Guessing \emph{any other} $k$ makes one branch $k-1$ or $m-k$ candidates -- both can be $>m/2$ --
so no single guess has a smaller worst case. That is the \propref{c25-halve} lower bound met exactly.
\end{tracebox}

\subsection{Variations and Twists}

\begin{itemize}[nosep]
  \item \textbf{Cost-aware split (LC 375)}: you \emph{pay} the guessed number and want to minimise
        worst-case total cash; optimal first guess is not the middle -- interval DP over
        $[i,j]$. \hfill\lv{3}--\lv{4}
  \item \textbf{Reconstruct a hidden \emph{tree} by distance queries}: query $\dist(u,v)$ against an
        oracle (answerable via \chref{c24}); adaptively locate a node or edge by which split shrinks
        the candidate-tree set -- group testing on a tree. \hfill\lv{4}--\lv{5}
  \item \textbf{Ancestor/yes/no queries}: ``is $u$ an ancestor of $v$?'' answered in $\bigO{1}$ by
        $\tin/\tout$ (\thmref{c24-anc}) is a binary query; $n$ of them rebuild parent pointers by
        bisecting each node's ancestor chain. \hfill\lv{4}
  \item \textbf{Non-adaptive (batch) queries}: if all questions must be fixed upfront (no adaptation),
        you need combinatorial group designs -- more queries than the adaptive $\log$. \hfill\lv{5}
  \item \textbf{Lies allowed}: Ulam's game -- each answer might be wrong $e$ times; the candidate
        count is spread over $\sum_{i\le e}\binom{k}{i}$ transcripts, inflating queries. \hfill\lv{5}
\end{itemize}

\subsection{Common Pitfalls}\label{sec:c25-query-pitfalls}

\begin{warningbox}[Trap 1 -- \code{(lo+hi)/2} overflowing]
With $n$ up to $2^{31}{-}1$, \code{lo+hi} exceeds \code{INT\_MAX} and wraps negative, corrupting the
midpoint (wrong answer or an infinite loop). Always write \code{lo + (hi-lo)/2} -- mathematically
equal, overflow-free.
\end{warningbox}
\begin{lstlisting}[caption={\textcolor{warningred}{WRONG}: overflows when \code{lo+hi} exceeds \code{INT\_MAX}}, language=C++, basicstyle=\ttfamily\small]
int mid = (lo + hi) / 2;         // lo,hi up to 2^31-1  ->  signed overflow (UB)
\end{lstlisting}
\begin{lstlisting}[caption={\textcolor{successgreen}{RIGHT}: overflow-safe midpoint}, language=C++, basicstyle=\ttfamily\small]
int mid = lo + (hi - lo) / 2;    // hi-lo can't overflow; result identical
\end{lstlisting}

\begin{warningbox}[Trap 2 -- flipping the oracle's sign]
\code{guess(mid)}$=-1$ means \emph{your guess is too high}, so the secret is \emph{smaller}: shrink
$\mathit{hi}$, not $\mathit{lo}$. Swapping the two branches makes the interval grow and the loop
never ends (or returns wrong). Read the oracle's contract literally.
\end{warningbox}

\begin{warningbox}[Trap 3 -- non-shrinking bounds $\Rightarrow$ infinite loop]
Updating to $\mathit{hi}=\mathit{mid}$ (instead of $\mathit{mid}-1$) while $\mathit{lo}$ stays can keep
$\mathit{mid}$ constant forever. Because you already know $\mathit{mid}$ is \emph{not} the answer on
that branch, exclude it: $\mathit{hi}=\mathit{mid}-1$ / $\mathit{lo}=\mathit{mid}+1$.
\end{warningbox}

\subsection{Problem Lab}

LC 374 is the section's cleanest specimen: a hidden integer, a three-way oracle, and the provably
optimal halving strategy -- with one real trap (overflow) hiding behind an ``easy'' label.

% ---------------------------------------------------------------------------
\solvedproblem{c25-guess}{Guess Number Higher or Lower}{LeetCode 374}
  {https://leetcode.com/problems/guess-number-higher-or-lower/}{LC Easy L2}
  {binary search, interactive/oracle model, three-way comparison, overflow-safe midpoint}

\begin{problemstatement}[Guess Number Higher or Lower]
We pick a number $x$ in $[1, n]$. You guess a number and a pre-defined function \code{guess(num)}
returns \textbf{$-1$} if your guess is \emph{higher} than $x$ (``too high''), \textbf{$1$} if your guess
is \emph{lower} than $x$ (``too low''), and \textbf{$0$} if correct. Return the number $x$.
$1\le n\le 2^{31}-1$.
\end{problemstatement}

\begin{attemptgate}[Attempt it yourself first -- \lv{2}, give yourself 12 minutes]
Fair first attempt: obviously binary search. But $n$ is enormous -- make sure your loop terminates
and your midpoint never overflows.
\peek{What is the largest value \code{lo+hi} can take when $n=2^{31}-1$, and what does \code{int}
do with it?}
\nudge{Keep a $[\mathit{lo},\mathit{hi}]$ window. Compute \code{mid} as \code{lo+(hi-lo)/2}. On
\code{guess(mid)}$=-1$ move \code{hi} below \code{mid}; on $1$ move \code{lo} above \code{mid}; on $0$
return.}
\fullhint{Standard three-way binary search with an overflow-safe midpoint and correct branch
direction (\code{$-1$ $\Rightarrow$ secret smaller}). Loop while \code{lo<=hi}. Full build and proof
below.}
\end{attemptgate}

\solstep{Understanding the Problem}
A monotone predicate: define $P(k)\equiv(\mathit{guess}(k)\le 0)$, i.e.\ $k\ge x$. Truth grows with
$k$, so $x$ is the smallest $k$ with $P(k)$ true -- exactly binary search on a sorted yes/no boundary.
The oracle gives three-way feedback ($<,=,>$), so you can finish the moment $\mathit{guess}(\mathit{mid})=0$.
The only genuine difficulty is numeric: $n$ reaches \code{INT\_MAX}, so careless midpoint arithmetic
overflows.

\solstep{Key Observations}
\begin{observation}
Each answer collapses the candidate interval to at most half if you guess the midpoint
(\propref{c25-halve}), so $\le\lceil\log_2(n{+}1)\rceil$ guesses (about $31$ for the largest $n$) always
suffice -- and no strategy guarantees fewer (\thmref{c25-logn}).
\end{observation}
\begin{observation}
\code{guess(mid)}$=-1$ means the secret is \emph{strictly below} \code{mid}; $=1$ means \emph{strictly
above}. Because \code{mid} is known wrong on that branch, exclude it (\code{mid$\pm$1}) -- this both
prevents infinite loops and keeps the count at $\log$.
\end{observation}
\begin{observation}
\code{lo+hi} can overflow \code{int}; \code{lo+(hi-lo)/2} cannot, since $0\le\mathit{hi}-\mathit{lo}
\le n-1$ fits. This one line is the difference between AC and a Wrong Answer on large $n$.
\end{observation}

\solstep{Think It Through With Me}
\begin{thinkbox}
\deadend{Linear scan $1,2,\dots$ until \code{guess}$=0$.}{$O(n)$ oracle calls -- at $n=2^{31}$ it is
hopeless (TLE) even though each call is cheap. Queries are the currency here.}
\deadend{Guess \code{(lo+hi)/2}.}{On the biggest inputs \code{lo+hi} overflows to a negative \code{int},
\code{mid} lands outside $[\mathit{lo},\mathit{hi}]$, and you loop or return wrong. The idea is right;
the \emph{arithmetic} betrays it.}
\aha{Binary search with \code{mid=lo+(hi-lo)/2}: halve the window each answer, move to
\code{mid$\pm$1} (mid is known-wrong on that branch), stop when \code{guess=0}. Optimal and overflow-free.}
\end{thinkbox}

\solstep{Solution}
Maintain a closed candidate interval $[\mathit{lo},\mathit{hi}]=[1,n]$. While non-empty, compute the
overflow-safe midpoint, ask \code{guess}: on $0$ return it; on $-1$ set $\mathit{hi}=\mathit{mid}-1$
(secret lower); on $1$ set $\mathit{lo}=\mathit{mid}+1$ (secret higher). Monotonicity guarantees the
secret stays inside the window; halving guarantees termination in $\log$ steps.

\solstep{Proof of Correctness}
\begin{proofbox}
\emph{Invariant}: the secret $x\in[\mathit{lo},\mathit{hi}]$ at every loop start (true initially). If
$\mathit{guess}(\mathit{mid})=-1$ then $x<\mathit{mid}$, so setting $\mathit{hi}=\mathit{mid}-1$ keeps
$x$ in the window; the $1$ case is symmetric; the $0$ case returns $\mathit{mid}=x$. Thus the invariant
is preserved and the returned value equals $x$. \emph{Termination}: each iteration shrinks
$\mathit{hi}-\mathit{lo}$ to at most $\lceil(\mathit{hi}-\mathit{lo})/2\rceil-1<\mathit{hi}-\mathit{lo}$,
so after $\le\lceil\log_2(n{+}1)\rceil$ steps the window empties (and must have hit $0$ first, since
$x$ is always inside it). \emph{Optimality}: by \propref{c25-halve}, any strategy needs
$\Omega(\log n)$ three-way answers to distinguish $n$ ordered secrets, and \thmref{c25-logn} shows
this attains the bound, so it is query-optimal. \end{proofbox}

\solstep{Implementation}
See \coderef{c25-guess} above -- the listing in this section \emph{is} the accepted solution:
three-way binary search with \code{lo+(hi-lo)/2} and correct branch directions.

\solstep{Trace}
\begin{tracebox}[Official $n=10, x=6$]
$[\mathit{lo},\mathit{hi}]=[1,10]$: $\mathit{mid}=5$, $\mathit{guess}=1$ (low) $\Rightarrow\mathit{lo}=6$.
$[6,10]$: $\mathit{mid}=8$, $\mathit{guess}=-1$ (high) $\Rightarrow\mathit{hi}=7$.
$[6,7]$: $\mathit{mid}=6$, $\mathit{guess}=0$ $\Rightarrow$ return $6$. \checkmark\ Three guesses for
$x=6$ -- within $\lceil\log_2 11\rceil=4$.
\end{tracebox}

\solstep{Complexity}
\begin{complexitybox}
Time $\bigO{\log n}$ oracle calls ($\le 31$ at the maximum $n$); each call is treated as $O(1)$.
Space $\bigO{1}$. Optimal by \thmref{c25-logn}. Verdict: \ac.
\end{complexitybox}

\solstep{Pitfalls and Variants}
\begin{warningbox}
(i) \code{(lo+hi)/2} overflow (Trap~1) is the single most common WA/TLE here. (ii) Sign flip of the
oracle (Trap~2). (iii) Using $\mathit{hi}=\mathit{mid}$ / $\mathit{lo}=\mathit{mid}$ instead of
$\pm1$ (Trap~3) -- infinite loop on the boundary. \emph{Variant:} LC 374's ``guess the word'' (1095)
and ``find peak element'' (162) are the same adaptive halving on a different oracle -- binary search
without an explicit $=$ answer.
\end{warningbox}

\begin{insightbox}[Key takeaways]
\begin{enumerate}[nosep, leftmargin=1.6em]
  \item \textbf{Choose questions that halve.} With an oracle, your power is picking the query whose
        answers split the candidate space evenly (\propref{c25-halve}).
  \item \textbf{Counting bound $=$ decision tree.} $\log$ questions is a tree of depth $\log$ -- the
        same object as every other chapter, but nodes are \emph{queries}.
  \item \textbf{Easy label, real trap.} ``Easy'' hid an integer-overflow bug; correctness of the
        \emph{arithmetic} matters as much as the algorithm.
\end{enumerate}
\end{insightbox}

\begin{hardeningbox}
\begin{itemize}[nosep, leftmargin=1.4em]
  \item \textbf{Prove optimality}: show no decision tree with fewer than $\lceil\log_2(n{+}1)\rceil$
        leaves covers $n$ ordered secrets -- a counting argument (\propref{c25-halve}).
  \item \textbf{Cost-aware version}: replace the midpoint with the min-max choice
        $\arg\min_k\big(k+\max(T[\mathit{lo},k{-}1],T[k{+}1,\mathit{hi}])\big)$ -- LC 375.
  \item \textbf{Tree reconstruction by distance}: to find a hidden node $x$, repeatedly query
        $\dist(\mathit{centroid},x)$ (answerable in $O(1)$ after \chref{c24}) to discard whole
        subtrees -- centroid-halving instead of line-halving (\chref{c09} if centroid known).
\end{itemize}
\end{hardeningbox}

\subsection{Practice Problems}

\begin{stuckbox}
\begin{itemize}[nosep, leftmargin=1.4em]
  \item Did you write \code{lo+(hi-lo)/2}, not \code{(lo+hi)/2}? (Trap~1.)
  \item Did you map \code{guess}$=-1$ to ``search \emph{lower}'' (shrink \code{hi})? (Trap~2.)
  \item Did you exclude the known-wrong \code{mid} with $\pm1$? (Trap~3.)
  \item Can you state why fewer than $\log_2 n$ guesses cannot always work? (\propref{c25-halve}.)
\end{itemize}
\end{stuckbox}

\begin{practiceset}
  \practice{Guess Number Higher or Lower II}{LeetCode 375}
    {https://leetcode.com/problems/guess-number-higher-or-lower-ii/}{LC Medium L3}
    {Same three-way oracle as \probref{c25-guess}, but a wrong guess of $k$ \emph{costs you \$k} and
    you must return the minimum cash needed to \textbf{guarantee} a win regardless of the secret. The
    halving midpoint is no longer optimal -- you now minimise the worst-case \emph{total cost}, a
    different objective on the same decision tree. Model it with interval DP: let $T[i][j]$ be the
    guarantee for secrets $i..j$; then $T[i][j]=\min_{k\in[i,j]}\big(k+\max(T[i][k-1],\,T[k+1][j])\big)$
    with $T[i][j]=0$ when $i\ge j$ (the adversary picks the pricier branch, so you prepare for the
    $\max$). This is the ``cost-aware split'' variation from \secref{c25-query}; it also generalises
    the optimal-BST / matrix-chain family -- an interval-DP technique, not a tree-representation one.}
\end{practiceset}

\begin{recallbox}[10-second refresher -- 6.4 Recover by Queries (Interactive)]
You only see the answer by \emph{asking}: a query model is (space, allowed questions, feedback). A
truthful answer is worth a branch, so $k$ $b$-way answers distinguish $b^{k}$ targets -- you cannot
beat $\lceil\log_b\rceil$ (\propref{c25-halve}). Guess the midpoint: three-way binary search finds a
secret in $[1,n]$ in $\lceil\log_2(n{+}1)\rceil$ guesses, provably optimal (\thmref{c25-logn}), with
\code{lo+(hi-lo)/2} to dodge \code{INT\_MAX} overflow and $\mathit{mid}\pm1$ to keep shrinking. When
wrong guesses \emph{cost}, minimise $k+\max(\text{left},\text{right})$ over the interval (LC 375); when
the hidden object is a tree, halve with distance/ancestor queries off the \chref{c24} oracle.
\selftest{
  \item State the query model triple and give the oracle for LC 374.
  \item Prove $\lceil\log_2(n{+}1)\rceil$ is a lower bound on worst-case guesses for a three-way oracle.
  \item Why does \code{(lo+hi)/2} fail at $n=2^{31}-1$, and what replaces it?
  \item In LC 375, why do we take $\max$ (not sum/average) of the two sub-intervals plus the cost $k$?
  \item Give one distance/ancestor query that would help reconstruct a hidden tree, and why it halves.
}
\end{recallbox}
